% Source-derived chapter 10: Questions and examples
% Original source: canonical-representations-surface-groups
\chapter{Questions and examples}
\label{canonical-representations-surface-groups-chapter-10}
\source{canonical-representations-surface-groups}

\begin{remark}[Source section]
\label{canonical-representations-surface-groups-chapter-10-source}
\source{canonical-representations-surface-groups}
\source{canonical-representations-surface-groups}
This chapter reproduces the corresponding section of the retrieved
LaTeX source, including its definitions, statements, examples, and
proofs. It has not yet been translated into Lean.
\notready
\end{remark}

% The source section title was: Questions and examples
\label{section:questions}
\source{canonical-representations-surface-groups}
\subsection{Rigidity questions for mapping class groups}
One of the key ingredients of our arguments is the statement that certain representations of (finite index subgroups of) $\on{Mod}_{g,n}$ are rigid.
\begin{question}
\source{canonical-representations-surface-groups}
\notready\label{question:rigidity}
\source{canonical-representations-surface-groups}
Let $g\gg 0$. Is every irreducible complex representation of every finite index subgroup of $\on{Mod}_{g,n}$ rigid?
\end{question}
The answer is, perhaps, plausibly ``yes" if one accepts the well-known analogy
between $\on{Mod}_{g,n}$ and lattices in simple Lie groups of rank $>1$, as
all complex representations of such lattices are rigid.  
For $g$ small, there
are examples of irreducible non-rigid representations of $\on{Mod}_{g,n}$ (for
example, $\pi_1(\mathscr{M}_{0,4})$ is free on two generators). 

\begin{remark}
\source{canonical-representations-surface-groups}
\notready
	\label{remark:}
\source{canonical-representations-surface-groups}
	A positive answer to \autoref{question:rigidity} would imply Ivanov's
conjecture that $\on{Mod}_{g,n}$ does not virtually surject onto $\mathbb{Z}$
as stated in \cite[\S 7]{Ivanov:problems} and \cite[Problem
2.11.A]{Kirby:problems}. Indeed, suppose some finite index subgroup
$\Gamma\subset \on{Mod}_{g,n}$ admitted a surjection $f: \Gamma
\twoheadrightarrow\mathbb{Z}$. For $t\in \mathbb{C}^\times$, the representation
\begin{align*}
\rho_t: \mathbb{Z} &\to \mathbb{C}^\times \\
n & \mapsto t^n
\end{align*}
is a non-trivial
family of one-dimensional representations of $\mathbb{Z}$. Then the
representation $\rho_t\circ f$ is a non-constant family of irreducible
representations of $\Gamma$.
\end{remark}

A positive answer to \autoref{question:rigidity} would also imply a positive answer to the following question:
\begin{question}
\source{canonical-representations-surface-groups}
\notready\label{question:finite-fixed-points}
\source{canonical-representations-surface-groups}
Let $g\gg0$. Let $\Gamma\subset \on{Mod}_{g,n}$ be a finite index subgroup, and
let $X_r(\pi_1(\Sigma_{g,n}))$ be the character variety parametrizing
$r$-dimensional semisimple representations of $\pi_1(\Sigma_{g,n})$. Is the fixed locus  $X_r(\pi_1(\Sigma_{g,n}))^\Gamma$ finite?
\end{question}

\begin{remark}
\source{canonical-representations-surface-groups}
\notready
One reason to believe \autoref{question:finite-fixed-points} has a positive answer is that the profinite analogue does. Namely, let $X$ be any normal connected algebraic variety over $\mathbb{C}$, and let $\pi_1^{\text{\'et}}(X)$ be its profinite fundamental group. Let $\Gamma$ be a finite index subgroup of $\on{Out}(\pi_1^{\text{\'et}}(X))$. Then $\Gamma$ acts on the set of conjugacy classes of continuous representations $$\pi_1^{\text{\'et}}(X)\to \on{GL}_r(\overline{\mathbb{Q}}_\ell).$$ It follows from \cite[Remark 1.1.4]{litt2021arithmetic} that the $\Gamma$-fixed points are \emph{discrete} in the $\ell$-adic topology. (The result there is stated for curves, but one may reduce to this case by the Lefschetz hyperplane theorem.)
\end{remark}



If one accepts in addition Simpson's motivicity conjecture for rigid local systems \cite[Conjecture on p.~9]{simpson1992higgs}, a positive answer to \autoref{question:rigidity} would imply that every semisimple local system on $\mathscr{M}_{g,n}$ is of geometric origin for $g\gg 0$. This would be of particular interest for those local systems not obviously of geometric origin, e.g.~the local systems arising from TQFT constructions. 
Thus one might ask:

\begin{question}
\source{canonical-representations-surface-groups}
\notready
For $g\gg 0$, is every semisimple local system on $\mathscr{M}_{g,n}$ of geometric origin? Let $C$ be a Riemann surface of genus $g\gg 0$. Is every semisimple MCG-finite representation of $\pi_1(C)$ of geometric origin?
\end{question}

\begin{example}
\source{canonical-representations-surface-groups}
\notready\label{non-rigid-examples}
\source{canonical-representations-surface-groups}
The analogy between the representation theory of mapping class groups and the representation theory of lattices in simple Lie groups of rank greater than one only goes so far. Indeed, mapping class groups typically admit non-rigid reducible representations, as we now explain. 
\begin{enumerate}
\item Let $H=H_1(\Sigma_g, \mathbb{Q})$. Morita \cite{morita1993extension} produces a nonzero class $$\sigma\in H^1(\on{Mod}_g, (\wedge^3 H)/H),$$ closely related to the Johnson homomorphism, yielding a non-split extension of $\on{Mod}_g$-representations $$1\to \wedge^3H/H\to W\to \mathbb{Q}\to 1.$$ The representation $W$ is evidently a non-trivial deformation of $\mathbb{Q}\oplus \wedge^3H/H$, and hence this latter representation is non-rigid.

\item An essentially identical but perhaps slightly simpler argument gives examples for punctured surfaces. For $n>0$ the group $\on{PMod}_{g,n}$ acts on $\pi_1(\Sigma_{g,n-1}, x)$, and hence on the group algebra $\mathbb{Q}[\pi_1(\Sigma_{g,n-1}, x)]$. Let $\mathscr{I}\subset \mathbb{Q}[\pi_1(\Sigma_{g,n-1}, x)]$ be the augmentation ideal. Then direct computation shows that the short exact sequence of $\on{PMod}_{g,n}$-representations $$0\to \mathscr{I}^2/\mathscr{I}^3\to \mathscr{I}/\mathscr{I}^3\to \mathscr{I}/\mathscr{I}^2\to 0$$ does not split, and hence the representation $\mathscr{I}/\mathscr{I}^2\oplus \mathscr{I}^2/\mathscr{I}^3$ is not rigid.

\item \label{h1-example} For another example, one may consider the non-torsion class in $$H^1(\on{Mod}_{g,1}, H^1(\Sigma_g, \mathbb{Z}))$$ constructed in \cite[Proposition 6.4]{morita1989families}, which yields a non-split extension of $\on{Mod}_{g,1}$-representations $$1\to H^1(\Sigma_g, \mathbb{Q})\to W\to \mathbb{Q}\to 1.$$ One may make the restriction to the point-pushing subgroup $W|_{\pi_1(\Sigma_g)}$ explicit as follows. Let $$a: \pi_1(\Sigma_g)\to \pi_1(\Sigma_g)^{\on{ab}}\otimes\mathbb{Q}\simeq H_1(\Sigma_g, \mathbb{Q})$$  be the map induced by the Hurewicz isomorphism $H_1(\Sigma_g, \mathbb{Z})=\pi_1(\Sigma_g)^{\on{ab}}$. Choose a splitting $W=H_1(\Sigma_g, \mathbb{Q})\oplus \mathbb{Q}e$, where we view $H_1(\Sigma_g, \mathbb{Q})$ as isomorphic to $H^1(\Sigma_g, \mathbb{Q})$ via the universal coefficient theorem and Poincar\'e duality. Then $\pi_1(\Sigma_g)$ acts trivially on $H_1(\Sigma_g, \mathbb{Q})$, and an element $\gamma\in\pi_1(\Sigma_g)$ acts on $e$ via the map $$e\mapsto e+a(\gamma).$$  
\source{canonical-representations-surface-groups}
\end{enumerate}
\end{example}

\begin{remark}
\source{canonical-representations-surface-groups}
\notready\label{remark:FLM-rank-1-result}
\source{canonical-representations-surface-groups}
	For a related discussion of ways mapping class groups \emph{do not} behave like lattices in simple Lie groups of rank greater than one, see \cite{farbLM:rank-1-phenomena}. Interestingly, \cite[Theorem 1.6]{farbLM:rank-1-phenomena} shows there are no faithful linear
representations of finite index subgroups of $\on{Mod}_{g,0}$ of dimension $< 2 \sqrt{g-1}$. A related result follows from our work here, as we now explain.

Let $\Gamma\subset \on{PMod}_{g,n+1}$ be a finite-index subgroup containing the point-pushing subgroup. 
It follows from \autoref{theorem:finite-image},
\autoref{lemma:finite-rho-implies-finite-lift}(1), and
\autoref{proposition:MCG-finite-index-rep}
that if $\rho$ is a representation of $\Gamma$ of rank $<\sqrt{g+1}$ 
\begin{enumerate}
\item 	whose restriction to the point-pushing subgroup is irreducible, and
\item with finite determinant
\end{enumerate}
then $\rho$ has finite image.
Although the two results are not comparable,
it is interesting that the bound in
\cite{farbLM:rank-1-phenomena} is also asymptotic to $\sqrt{g}$. 
\end{remark}

\subsection{Bounds and examples}\label{question:bounds}
\source{canonical-representations-surface-groups}
It is natural to ask how sharp the bound of $\sqrt{g+1}$ in
\autoref{theorem:finite-image} is. We have no reason to believe it is sharp.
That said, one cannot expect a bound that is much stronger, as the extension
$W$ constructed in \autoref{non-rigid-examples}(\ref{h1-example}) yields a
non-trivial unipotent representation of $\pi_1(\Sigma_g)$ of rank $2g+1$, fixed
by the action of $\on{Mod}_g$, namely $W|_{\pi_1(\Sigma_g)}$.  

\begin{remark}
\source{canonical-representations-surface-groups}
\notready
	\label{remark:sharp-g-1}
\source{canonical-representations-surface-groups}
	The bounds of many of our main results are sharp when $g=1$. 
In $g = 1$, there are semisimple $2$-dimensional ``special dihedral" MCG-finite representations
of $\pi_1(\Sigma_{1,n})$ for $n > 0$
\cite[Theorem B]{biswas2017surface},
with image contained in the infinite dihedral group.
This shows sharpness of \autoref{theorem:finite-image} and also the semisimple
case, \autoref{theorem:finite-image-semisimple}.
Moreover, there exist $2$-dimensional non-semisimple MCG-finite representations of $\pi_1(\Sigma_{1,n})$ for $n>0$, see \cite[Theorem B]{CH:isomonodromic}.
This implies sharpness of the bound in \autoref{lemma:semisimplicity}.

Note that the ``special dihedral" MCG-finite representations of \cite[Theorem B]{biswas2017surface} are not in general arithmetic in the sense of \autoref{definition:arithmetic}, as in general their local monodromy at infinity is not quasi-unipotent.
\end{remark}
\begin{remark}
\source{canonical-representations-surface-groups}
\notready
	\label{remark:not-sharp-g-2-3}
\source{canonical-representations-surface-groups}
	On the other hand, the bound $r < \sqrt{g+1}$ in \autoref{theorem:finite-image} is not sharp when $g =
	2$ and $g =3$.
	Indeed, to verify this, we only need show there are no $2$-dimensional
	MCG-finite representations with infinite image. There are no non-semisimple such
	representations by \cite[Theorem B]{CH:isomonodromic}, and so it only
	remains to show there are no irreducible $2$-dimensional
	representations.
	Any such representation $\rho: \pi_1(\Sigma_{g,n}) \to \on{GL}_2(\mathbb
	C)$ has finite determinant by \autoref{lemma:finite-central-part}.
	Replacing $\rho$ by $\rho \otimes \det \rho^{-1/2}$, we would obtain another
	MCG-finite $2$-dimensional representation with infinite image factoring
	through $\on{SL}_2(\mathbb C)$.
	No such representations	exist by \cite[Theorem A]{biswas2017surface}.
\end{remark}


\begin{question}
\source{canonical-representations-surface-groups}
\notready
Let $g, n\geq 0$. What is the minimal rank of an MCG-finite representation of $\pi_1(\Sigma_{g,n})$ with infinite image?
\end{question}

	It is plausible that much stronger bounds 
than $\sqrt{g+1}$ hold in \autoref{theorem:finite-image}
if one assumes the representation is semisimple (see \autoref{figure:MCG-geography}).

\begin{example}
\source{canonical-representations-surface-groups}
\notready
	\label{example:semisimple-bound}
\source{canonical-representations-surface-groups}
We know of examples of MCG-finite
semisimple representations of $\pi_1(\Sigma_g)$ with infinite image, of rank exponential in $g$. For example, one may via TQFT techniques construct representations of $\on{Mod}_{g,n}$
of rank exponential in $g$
\cite[Corollary 4.3]{koberda2016quotients}, \cite[Theorem 5.1]{biswas2018representations}; 
these representations are non-trivial when restricted to the point-pushing subgroup, and hence by \autoref{proposition:MCG-finite-index-rep} yield MCG-finite representations.
Similarly, representations of $\on{Mod}_{g,n}$ constructed via variants of
the Kodaira-Parshin trick 
\cite[Example 3.3.1]{lawrence2019representations}
(see also \cite[Proposition 5.1.1 and Remark 5.1.2]{LL:geometric-local-systems})
are semisimple of rank exponential in $g$ and restrict to 
MCG-finite, semisimple representations of the point-pushing subgroup, with infinite image. 
\end{example}

\begin{example}
\source{canonical-representations-surface-groups}
\notready
	In genus zero, there is a huge collection of MCG-finite representations, namely the \emph{rigid local systems} studied by Katz \cite{katz2016rigid}. The use of rigidity in our proof of \autoref{theorem:finite-image}, and \autoref{question:rigidity}, suggest that the study of MCG-finite representations is a natural generalization  of the study of rigid local systems to the higher genus setting. 
\end{example}

\begin{question}
\source{canonical-representations-surface-groups}
\notready
Let $g, n\geq 0$. What is the minimal rank of a semisimple MCG-finite representation of $\pi_1(\Sigma_{g,n})$ with infinite image?
\end{question}

In our view it would be extremely interesting to give an improved bound for semisimple MCG-finite representations, and to produce fundamentally new examples of semisimple MCG-finite representations.

\begin{remark}[Bounds for free groups]
\source{canonical-representations-surface-groups}
\notready
	\label{example:free-group-bounds}
\source{canonical-representations-surface-groups}
	A variant of the construction in \autoref{non-rigid-examples}(\ref{h1-example})
gives an example of a representation of the free group $F_N$ on $N$ generators
of rank $N+1$, whose conjugacy class is fixed by the action of $\on{Out}(F_N)$,
see \cite[p.~1444]{potapchik2000low}. For example, when $N=2$, this representation is given by the matrices $$\begin{pmatrix} 1 & 0 & 1 \\ 0 & 1 & 0 \\ 0 & 0 & 1 \end{pmatrix}, \begin{pmatrix} 1 & 0 & 0 \\ 0 & 1 & 1 \\ 0 & 0 & 1 \end{pmatrix}.$$
Thus again the bound of $\sqrt{g+1}$ in
\autoref{corollary:free-groups} cannot be improved too much further if one
allows non-semisimple representations. However it is in principle possible that
there is \emph{no} semisimple representation of $F_N$ with infinite image whose
conjugacy class has finite orbit under $\on{Out}(F_N)$ when $N\geq 3$. Indeed, this would
follow from a conjecture of Grunewald and Lubotzky, see \cite[\S9.2]{grunewald2009linear}, using the main result of \cite{farb2017moving}.

There are interesting semisimple representations of $F_2$ whose conjugacy class
has finite orbit under $\on{Out}(F_2)$. This follows, for example, from the linearity of $\on{Aut}(F_2)$ \cite{krammer2000braid}, e.g.~by choosing a faithful representation of $\on{Aut}(F_2)$ and restricting to the inner automorphisms, and then semisimplifying. Alternately, one may construct examples by applying the Kodaira-Parshin trick to construct representations of $\on{Mod}_{1,2}$, which is commensurable with $\on{Aut}(F_2)$.
%We now explain why the above claim follows from
%\cite[\S9.2]{grunewald2009linear}, which conjectures every representation
%$\xi : \on{Aut}(F_N) \to \on{GL}(V)$ of
%$\on{Aut}(F_N)$ restricts to a representation on $\on{Inn}(F_N)\simeq F_N$ with virtually
%solvable image. 
%Suppose we are given a semisimple representation $\rho$ of $F_N$ with finite orbit under
%$\on{Out}(F_N)$.
%In order to show $\rho$ has finite image, we may assume $\rho$ is irreducible.
%Then, using Schur's lemma as in \autoref{lemma:mcg-rep-construction},
%there is a finite index subgroup 
%$\Gamma \subset \on{Aut}(F_N)$
%containing $\on{Inn}(F_N)$
%and an irreducible representation $\xi: \Gamma \to \on{PGL}_r(\mathbb C)$
%restricting to the projectivization of $\rho$ on $\on{Inn}(F_N) \simeq F_N$.
%Composing $\xi$ with an embedding 
%$\on{PGL}_r \to \on{GL}_{r'}$, 
%and then inducing this from $\Gamma$ to $\on{Aut}(F_N)$, we obtain a
%semisimple representation $\xi'$ of $\on{Aut}(F_N)$ which restricts to
%a semisimple representation of $F_N$.
%Assuming the conjecture \cite[\S9.2]{grunewald2009linear},
%the Zariski closure of $\on{Im}(\xi)$ has identity component which is a solvable semisimple
%group.
%Therefore, it is a torus. Hence, there is a finite index subgroup of 
%$\on{Im}(\rho)$ whose projectivization is contained in a torus.
%This means moreover that there is a finite index subgroup $K \subset \on{Im}(\rho)$
%so that
%$\rho|_K$ has image which is Zariski dense in a torus $T$.
%Therefore, $\rho|_K$ has abelian image and hence factors through the
%abelianization of $K$, i.e., $H_1(K)$.
%By shrinking $K$, we may assume $K$ is characteristic in $F_N$.
%By \cite{farb2017moving}, since $N \geq 3$, the action of $\on{Out}(F_N)$
%on $H_1(K)$ has no nonzero vectors with finite orbit.
%On the contrary, the action of $\on{Out}(F_N)$ preserves the torus $T$, and
%hence acts on $T$ via its finite order Weyl group.
%Therefore, any infinite order element of $\rho(K)$, which must exist by finite
%generation of $K$ if $T$ is nontrivial, is fixed by a finite index
%subgroup of $\on{Out}(F_N)$, of index at most that of the Weyl group.
%This contradicts \cite{farb2017moving} above.
%This implies $T$ must be trivial, so $\on{Im}(\rho)$ is finite.
\end{remark}

As remarked in the introduction, we conjecture the opposite of Grunewald and Lubotzky:
\begin{conjecture}
\source{canonical-representations-surface-groups}
\notready
	\label{conjecture:infinite-image-semisimple}
\source{canonical-representations-surface-groups}
For all $N>0$, there exist semisimple representations of $F_N$ with infinite image, whose conjugacy class has finite orbit under $\on{Out}(F_N)$. 
\end{conjecture}
Our main impetus for this conjecture is the existence of many interesting
MCG-finite representations of $\pi_1(\Sigma_{g,n})$, see \autoref{example:semisimple-bound}. In our view this conjecture is of great importance, both because of the intrinsic interest of ``canonical" representations of $F_N$ (i.e.~those with finite orbit under $\on{Out}(F_N)$), and because of its relationship to the representation theory of $\on{Aut}(F_N)$. 

\begin{remark}
\source{canonical-representations-surface-groups}
\notready
We briefly spell out the the relationship between representations of $F_N$ whose conjugacy class has finite orbit under $\on{Out}(F_N)$, and the representation theory of
$\on{Aut}(F_N)$.
If $\Gamma\subset \on{Aut}(F_N)$ is a finite index subgroup containing the inner automorphisms, and
$$\rho:\Gamma\to \on{GL}_r(\mathbb{C})$$ is a representation, then the conjugacy
class of $\rho|_{F_N}$ has finite orbit under $\on{Out}(F_N)$, via a proof
analogous to that of
\autoref{proposition:MCG-finite-index-rep}. Conversely, as in to the proof of \autoref{lemma:mcg-rep-construction}, the projectivization of any irreducible $F_N$-representation whose conjugacy class has finite orbit under $\on{Out}(F_N)$ arises from a projective representation of a finite index subgroup of $\on{Aut}(F_N)$. See \cite[Theorem 7.12]{baumeister-kielak-pierro} for a related result about low-rank projective representations of certain finite index subgroups of $\on{Aut}(F_n)$.
\end{remark}
\subsection{Classification}
The following is evidently quite difficult, but in our view is of great interest:
\begin{question}
\source{canonical-representations-surface-groups}
\notready
Can one classify MCG-finite representations of $\pi_1(\Sigma_{g,n})$	 of rank $r$? Equivalently, can one classify algebraic solutions to the rank $r$ isomonodromy differential equations over $\mathscr{M}_{g,n}$? Can one classify representations of the free group $F_N$ on $N$ generators whose conjugacy class has finite orbit under $\on{Out}(F_N)$?
\end{question}
\begin{remark}
\source{canonical-representations-surface-groups}
\notready
	\label{remark:}
\source{canonical-representations-surface-groups}
It would also be extremely interesting to learn of any new sources of
\emph{examples} of MCG-finite representations. For example, are there MCG-finite representations $\pi_1(\Sigma_{g,n})\to \on{SL}_r(\mathbb{C})$ with (Zariski-)dense image for $r$ arbitrarily large?
\end{remark}


\appendix
